\documentclass{situationist-preprint}

\title{The Value of Knowledge and the Statues of Daedalus}
\msid{EPI-14-a}
\msfamily{EPI-14}
\msdiscipline{Epistemology}
\mstype{Principal manuscript}
\mskeywords{value of knowledge; Meno; swamping problem; stability; safety; reliabilism; misleading evidence}

\begin{document}
\maketitle

\begin{abstract}
In the \emph{Meno}, Socrates observes that true opinion is as good a guide to Larissa as knowledge, and then explains why knowledge is nonetheless prized more: true opinions, like the statues of Daedalus, run away unless they are tied down, and knowledge is true opinion tied down. We defend this answer. Knowledge is more valuable than mere true belief because it is more stable, and the stability is a source of value. The swamping problem holds that a property of a belief valued only as a means to truth adds nothing to a belief that is already true. We reply that stability is valued as a means to future true belief, which the present truth of the belief does not contain. We then explain why knowledge is more stable. Evidence that arrives after a belief is formed is produced by the facts; a safe belief, which could not easily have been false on its basis, is connected to the facts that produce later evidence, while a belief whose truth is lucky relative to its basis is exposed to evidence that its basis was defective. Under rational updating, mere true belief is therefore abandoned more often than knowledge. We answer the objections from stubbornness, from misleading evidence and from one-shot action.
\end{abstract}

\section{The problem}

Socrates asks Meno to consider a man who knows the road to Larissa and a man who has a correct opinion about it without knowing it. Each will lead others there correctly, and as far as getting to Larissa is concerned, true opinion is no worse a guide than knowledge \citep[97a--c]{PlatoMeno}. Why then is knowledge prized above true opinion? Socrates' answer is an image. The statues of Daedalus, if not tied down, run away; tied down, they stay. True opinions are fine things while they stay, but they do not stay long; they run away from the soul, and they are not worth much until one ties them down by working out the reason. When they are tied down they become knowledge, and then they are stable. That is why knowledge is more valuable than right opinion \citep[97d--98a]{PlatoMeno}.

The contemporary form of Meno's question is the \emph{swamping problem} \citep{Jones1997,Kvanvig2003,Zagzebski2003}. Suppose the only epistemic value is truth, and that the properties that make a true belief knowledge, such as being reliably formed, are valuable only as means to truth. Then once a belief is true, its being reliably formed adds nothing, just as a cup of excellent coffee is no better for having come from a reliable machine \citep{Zagzebski2003}. The value of the means is swamped by the value of the end. So, if truth is the only fundamental epistemic value, knowledge is not more valuable than true belief.

Responses divide. Some locate a final value in knowledge as an achievement creditable to the knower \citep{Greco2010,Sosa2007}. Some give up the claim that knowledge is distinctively valuable and locate the value elsewhere, in understanding \citep{Kvanvig2003} or in a weaker status \citep{PritchardMillarHaddock2010}. Some return to Plato: \citet[ch.~3]{Williamson2000} argued that knowledge is more robust than mere true belief under the impact of new evidence, and \citet{Olsson2007} and \citet{GoldmanOlsson2009} argued that a reliably produced true belief raises the probability of future true beliefs of the same kind. We defend the Platonic answer in a form that connects it with the analysis of knowledge as safe, justified, true belief (EPI-01), and that answers the swamping problem directly.

\section{Framework}

\begin{definition}[Stability]\label{d:stab}
A belief is more \emph{stable} than another if, under rational responses to the evidence that the believer subsequently receives, it is retained for longer in expectation.
\end{definition}

Stability so defined is not stubbornness. A stubborn belief is retained against evidence; a stable belief is retained in accordance with it, because the evidence does not tell against it.

\begin{principle}[Temporal value]\label{p:temporal}
The value of holding a true belief that $p$ accrues over the time during which it is held and available to guide action and inference.
\end{principle}

\ref{p:temporal} is Socrates' premise. The guide to Larissa is valuable while he guides; a true belief that runs away before the journey is over guides no further. It does not say that truth at a time has no value; it says that a true belief has value as a standing state, which persists and is used, and that its value as a state depends on how long it persists.

\begin{assumption}[Facts produce evidence]\label{a:produce}
Much of the evidence that bears on whether $p$, and that a believer receives after forming a belief about $p$, is produced by the facts that make $p$ true or false, or by the facts on which her basis for believing $p$ depends.
\end{assumption}

A traveller who believes that the road to Larissa is the left fork will see signposts, meet travellers coming from Larissa, and observe the landscape; all of this is produced by where the road goes. A traveller who believes it on the basis of a map will also, in time, receive evidence about the map: other maps, corrections, the map's errors elsewhere.

\section{The argument}

\begin{lemma}[Exposure]\label{l:exposure}
If a justified true belief is unsafe, there is a defect in its basis that the facts can reveal; if it is safe, there is none.
\end{lemma}

\begin{proof}
By EPI-01, a justified true belief is unsafe exactly when there are close cases in which the subject believes $p$ on the same basis and $p$ is false: its truth is lucky relative to its basis. The basis is then one that does not discriminate the actual case from those close cases, and this is a fact about the basis and the circumstances, the stopped clock, the barn facades, the misprinted map, which is itself among the facts of the world. A safe belief has a basis that delivers $p$ only in cases where $p$, and has no such defect with respect to $p$.
\end{proof}

\begin{theorem}[Stability]\label{t:stab}
Knowledge is more stable than mere justified true belief.
\end{theorem}

\begin{proof}
Consider a believer who knows $p$ and her duplicate who has a justified true belief in $p$ that is unsafe, and suppose both respond rationally to subsequent evidence. By Assumption~\ref{a:produce}, much of their subsequent evidence is produced by the facts. For the knower, the facts include $p$ and a basis without defect, so the evidence they produce tends to confirm $p$ and gives no reason to doubt her basis; only misleading evidence, which is not produced by the facts in this way, can lead her rationally to abandon her belief. For the duplicate, the facts include, by Lemma~\ref{l:exposure}, a defect in her basis; evidence of that defect is produced along with the rest, and when she receives it, rational updating requires her to withdraw her belief or to hold it only on a new basis. So the duplicate is exposed, in addition to misleading evidence, to a source of rational abandonment from which the knower is protected. Her belief is therefore retained for a shorter time in expectation, and the knower's belief is more stable.
\end{proof}

The argument can be made quantitative in the simplest case. Suppose that, in each period, the chance that a believer receives evidence that rationally leads her to give up a belief is $m$ if the evidence is misleading, and that an unsafe belief faces in addition a chance $d > 0$ of revelation of its defect. If these are independent from period to period, the expected lifetime of the safe belief is $1/m$ and of the unsafe one $1/(m + d)$, which is shorter by a factor of $m/(m + d)$. The man who has the right opinion about the road on the basis of a misread map will, on most journeys, reach a fork where the map's error shows.

\begin{corollary}[Value]\label{c:value}
Other things being equal, knowledge that $p$ is more valuable than mere true belief that $p$.
\end{corollary}

\begin{proof}
By Theorem~\ref{t:stab}, knowledge is retained for longer in expectation; by \ref{p:temporal}, the value of a true belief accrues over the time it is held. A mere true belief that is not justified is not more stable than an unsafe justified one, since its lack of justification is a further ground for rational abandonment; so the comparison extends from justified true belief to true belief in general.
\end{proof}

\paragraph{Swamping.} The swamping argument assumes that the extra value of knowledge must be the value of a means to the truth of the belief that $p$ at the time. Stability is a means to something else: the future truth-possession of the believer, her having a true belief about $p$ at later times. A belief that is true now does not contain that future good, and the value of the means is not swamped by a value already present. The coffee analogy fails at this point: the cup of coffee is consumed once, and the machine's reliability adds nothing to it; a belief is not consumed in use, and its value depends on its staying. This answer is close to that of \citet{GoldmanOlsson2009}, with one difference. Their extra value lies in the probability of future true beliefs on other questions, produced by the same reliable process; ours lies in the persistence of the same true belief. Both are means to future true belief, and the two answers are compatible.

\section{Objections}

\begin{objection}[Stubbornness]
A stubborn believer who never abandons her true beliefs has more stability than any knower. If stability made beliefs valuable, stubbornness would be an epistemic virtue.
\end{objection}

\begin{reply}
Definition~\ref{d:stab} measures stability under rational responses to evidence. The stubborn believer keeps her true beliefs, and keeps her false ones with them; her disposition, applied across her beliefs, produces no expected gain in truth and a large expected loss. The stability of knowledge is the stability of a belief that the evidence, rationally assessed, does not dislodge, and it is gained without any such cost. This is why Socrates says that true opinions are tied down \enquote{by working out the reason}, and not by refusing to listen.
\end{reply}

\begin{objection}[Misleading evidence]
Knowledge can be lost to misleading evidence. A person who knows that the dictator has been assassinated, having read it in a newspaper, may lose that knowledge when all other papers carry false denials \citep{Harman1973}. So knowledge is not stable.
\end{objection}

\begin{reply}
The claim is comparative. Misleading evidence threatens knowledge and mere true belief equally, since it does not discriminate between them; in the quantitative model it contributes $m$ to the hazard of both. What Theorem~\ref{t:stab} establishes is that mere true belief faces an additional hazard. Socrates' statues, tied down, can still be stolen; they do not walk away by themselves.
\end{reply}

\begin{objection}[One-shot action]
If a belief is used only once, for a single decision, its stability is irrelevant, and knowledge is no more valuable than true belief. Socrates' own case is of this kind: one journey to Larissa.
\end{objection}

\begin{reply}
We accept the consequence for the belief considered as an input to that decision: at the moment of action, true belief is as good a guide as knowledge, as Socrates says. The value claim concerns the belief as a state, which is the object of Meno's question and of ours. In any case, the journey to Larissa is not one-shot: the traveller must keep his belief through every fork, and a belief that runs away at the third fork is a worse guide than one that does not.
\end{reply}

\begin{objection}[Final value]
The interesting claim is that knowledge has more final value than true belief, and stability yields only instrumental value. The Platonic answer changes the subject \citep{PritchardMillarHaddock2010}.
\end{objection}

\begin{reply}
The question is whether knowledge is more valuable than true belief, and an instrumental difference is a difference in value. We do not deny that knowledge may also have final value as an achievement \citep{Greco2010}; the stability answer neither requires nor excludes it. It has the advantage of being available to an epistemology on which truth is the only fundamental epistemic value, which is the epistemology the swamping argument targets.
\end{reply}

\section{Conclusion}

Knowledge is more valuable than mere true belief because it is more stable. The stability is not a separate ingredient added to knowledge but a consequence of safety: a belief that could not easily have been false on its basis is connected to the facts that produce later evidence, and so survives rational updating that exposes a lucky belief's defective basis. The value of stability is the value of continuing to have a true belief, which the present truth of a belief does not contain, and so it is not swamped. Socrates' image of the tethered statues states the view exactly.

\begin{thebibliography}{99}\small
\bibitem[Goldman and Olsson(2009)]{GoldmanOlsson2009} Goldman, A.~I. and Olsson, E.~J. (2009). Reliabilism and the value of knowledge. In A.~Haddock, A.~Millar and D.~Pritchard (Eds.), \emph{Epistemic Value}. Oxford University Press.
\bibitem[Greco(2010)]{Greco2010} Greco, J. (2010). \emph{Achieving Knowledge: A Virtue-Theoretic Account of Epistemic Normativity}. Cambridge University Press.
\bibitem[Harman(1973)]{Harman1973} Harman, G. (1973). \emph{Thought}. Princeton University Press.
\bibitem[Jones(1997)]{Jones1997} Jones, W.~E. (1997). Why do we value knowledge? \emph{American Philosophical Quarterly}, 34(4), 423--439.
\bibitem[Kvanvig(2003)]{Kvanvig2003} Kvanvig, J.~L. (2003). \emph{The Value of Knowledge and the Pursuit of Understanding}. Cambridge University Press.
\bibitem[Olsson(2007)]{Olsson2007} Olsson, E.~J. (2007). Reliabilism, stability, and the value of knowledge. \emph{American Philosophical Quarterly}, 44(4), 343--355.
\bibitem[Plato(c.~385 BCE)]{PlatoMeno} Plato (c.~385 BCE). \emph{Meno}.
\bibitem[Pritchard et~al.(2010)]{PritchardMillarHaddock2010} Pritchard, D., Millar, A. and Haddock, A. (2010). \emph{The Nature and Value of Knowledge: Three Investigations}. Oxford University Press.
\bibitem[Sosa(2007)]{Sosa2007} Sosa, E. (2007). \emph{A Virtue Epistemology: Apt Belief and Reflective Knowledge, Volume I}. Oxford University Press.
\bibitem[Williamson(2000)]{Williamson2000} Williamson, T. (2000). \emph{Knowledge and Its Limits}. Oxford University Press.
\bibitem[Zagzebski(2003)]{Zagzebski2003} Zagzebski, L. (2003). The search for the source of epistemic good. \emph{Metaphilosophy}, 34(1--2), 12--28.
\end{thebibliography}

\end{document}
